\begin{papercard}{schaub2015}
\cf{Stages} \stagetag{S} \stagetag{C} (also \stagetag{W} \stagetag{H}). A systematisation-of-knowledge paper. It sets out requirements and best practices for privacy notices and a four-dimension \emph{design space}: timing, channel, modality and control. It is a ready-made coding frame for describing where and how an app communicates privacy.

\cf{Method of construction} A literature survey plus analysis of existing notices, built by design-science methodology (Peffers et al.). The initial taxonomy was refined with informal feedback from about 20 privacy experts and tested for expressiveness on scenarios \pc{schaub2015}{6}{Following Peffers et al.’s research methodology [102], we developed and refined the design space}. The options are not exclusive and the space can be extended \pc{schaub2015}{6}{the options for each di}.

\cf{Roles and hurdles} Notices serve users, companies (compliance, trust) and regulators (enforcement). Hurdles: \emph{notice complexity}, e.g.\ an estimated 244 hours a year to read the policies an American user encounters \pc{schaub2015}{2}{would take about 244 hours per year}; \emph{lack of choices} (take it or leave it); \emph{notice fatigue} and habituation \pc{schaub2015}{3}{Frequent exposure to seemingly irrelevant privacy notices results in habituation}; and \emph{decoupled notices}, e.g.\ fitness trackers with a policy only on the website \pc{schaub2015}{3}{many manufacturers of fitness tracking devices provide a privacy policy on their websites}.

\cf{Requirements and best practices} (1) Understand privacy in the system through a PIA: information flows, stakeholders, data minimisation \pc{schaub2015}{3}{The first step in designing effective privacy notices is to understand a system’s information flows}. (2) Different notices for different audiences: primary, secondary and incidental users \pc{schaub2015}{4}{Typical audience groups are the pri}. (3) Relevant, actionable information: prioritise \emph{unexpected} practices that cross contextual boundaries, i.e.\ violate contextual integrity \pc{schaub2015}{4}{Those practices are relevant because they cross contextual boundaries}. For each practice, gather who collects, why, use, retention, deletion and sharing recipients \pc{schaub2015}{4}{For each data practice, all parameters relevant for creat}. Add reasons, risks and remedies so that the notice is actionable \pc{schaub2015}{4}{This can be achieved by iden}. (4) System constraints and opportunities. (5) Layered, contextualised notices: the smallest layer must capture the main aspects or point to fuller notices \pc{schaub2015}{5}{the smallest notice either already captures the main aspects of the data practice}. Avoid habituation by prioritising notices by risk \pc{schaub2015}{5}{A good practice is to prioritize what and when notices are shown based on privacy risks}. (6) User-centred evaluation of attention, comprehension, recall and informed choice, or expert heuristic evaluation when resources are short \pc{schaub2015}{5}{In such cases, expert evaluation with usability heuristics}.

\cf{Design space (Figure 1)} \textbf{Timing}: at setup, just in time, context-dependent, periodic, persistent, on demand \pc{schaub2015}{6}{We describe six possible timing opportunities here}. Setup notices meet habituated users who cannot yet weigh trade-offs \pc{schaub2015}{6}{At setup time, users may have difficulty making informed decisions}. Just-in-time notices suit sensitive or unexpected practices such as location, contacts, photos and microphone access \pc{schaub2015}{6}{Just-in-time notices and obtaining express consent are particularly relevant}. On-demand: a persistent policy, or better a dashboard with summaries, controls and deletion, plus privacy-office contact details \pc{schaub2015}{8}{A better option are privacy settings interfaces or privacy dashboards}. \textbf{Channel}: primary (same device and context), secondary (out of band: website, e-mail, companion app), public (signs, beacons) \pc{schaub2015}{8}{the notice is provided within the user’s interaction with the system}. \textbf{Modality}: visual (text, icons, LEDs; wording, translation and personalisation matter), auditory, haptic, machine-readable \pc{schaub2015}{9}{Failing to translate may leave international users}. Icons such as AdChoices are poorly understood, so combine icons with text \pc{schaub2015}{9}{combining icons with textual explanations in privacy notices may improve}. \textbf{Control}: blocking (must decide to continue), non-blocking, decoupled (dashboards, settings menus). Defaults matter because they are read as recommendations \pc{schaub2015}{10}{as they may be kept unchanged out of convenience or because they are in}, \pc{schaub2015}{10}{Until the user provides a choice he or she cannot continue}.

\cf{Use cases} Websites mostly offer on-demand, primary, visual notices with decoupled controls \pc{schaub2015}{11}{Websites typically pro}. Smartphones in 2015: Android showed an all-or-nothing permission list at install (blocking, at setup); iOS asked just in time with per-resource control \pc{schaub2015}{12}{Thus, the main difference between the platforms is the level of control}. Lifelogging leaves bystanders without notice.

\cf{Conclusion} A privacy policy is necessary for compliance but insufficient for informing users. Notices must provide actionable choices \pc{schaub2015}{12}{A key aspect of effective notice design is the realization that a privacy policy}.

\cf{Use for this thesis} A theory-backed \emph{codebook dimension set} for privacy communication in apps. For each privacy-relevant item found in the walkthrough, code timing (store page, onboarding, just-in-time permission, settings on demand), channel (in app vs.\ website or e-mail), modality (text, icon, layered) and control (blocking, non-blocking, decoupled). The ``unexpected practices'' criterion links notice design to contextual integrity (Nissenbaum) for the harm stage. The audience list (primary user, partner, bystanders) matters for intimate health data. The 2015 platform details are outdated; check the current Android and iOS permission models before citing them.
\end{papercard}
